% ---------------------------------------------------------------------------
% INTRODUCTION --- v4 (2026-08-26). Register calibrated against DextrAH-G's
% actual introduction (arXiv:2407.02274v2), fetched verbatim.
%
% CHANGE FROM v3 (important, do not undo):
%   Stiffness is NO LONGER a pillar of the argument. The policy observes only a
%   point cloud, so it cannot be stiffness-ADAPTIVE by construction; leading on
%   stiffness would advertise a problem we do not solve. Our own failure data
%   implicates POSE and SIZE, not stiffness. Stiffness now appears once, in P4,
%   as honest ROBUSTNESS (material randomisation -> a grip that works across the
%   range, not one tuned to a value). Prepared answer if a reviewer presses:
%   "the policy is not stiffness-adaptive and cannot be; randomising material
%   during demonstration generation yields robustness across the range. Whether
%   stiffness-adaptive behaviour would help is open and we do not address it."
%
%   2026-08-26: contribution 2 further softened to DESCRIBE the training
%   ("across randomised object pose, size, shape, and material") rather than
%   CLAIM an outcome ("robust across the material range"). We have no stiffness
%   sweep demonstrating robustness; do not reinstate the outcome wording unless
%   that experiment exists.
%
%   Also restores the Fig. 1 reference (the 50-demo coverage-failure map), which
%   v3 dropped --- it is our own strongest motivating evidence.
%
% Structure: P1 problem | P2 prior work + what it leaves | P3 coverage evidence
% (cited) + our own failure map | P4 proposal (three beats only: FEM stress-aware
% scripted demonstrator -> many categories with within-category randomisation ->
% one policy, deployed real) | P5 contributions, standalone.
%
% P4 is deliberately NOT a full pipeline walkthrough (user, 2026-08-26). Keep it
% to the three beats; detail belongs in Method. The specialist/self-distillation
% stage is intentionally absent from the spine --- it is ablated, and leaving it
% out means the paragraph survives either outcome.
% ---------------------------------------------------------------------------
\newcommand{\method}{\textsc{Method}}   % TODO: name the framework

% Pending numbers --- defined once so a landing result is a one-line edit.
% Keep them bold while unresolved so they are impossible to miss in the PDF.
\newcommand{\ndemos}{\textbf{[XX\,k]}}      % synthesised demonstrations
\newcommand{\ncats}{\textbf{[N]}}           % fragile food categories
\newcommand{\nheld}{\textbf{[M]}}           % held-out categories (zero-shot), if reported
\newcommand{\nreal}{\textbf{[55]}}          % real teleop demonstrations in the co-train slice
\newcommand{\headline}{\textbf{[XX\%]}}     % headline real-robot success rate

\section{INTRODUCTION}

Handling delicate food is a skill people exercise without thought, yet it remains difficult for
robots. Automating it matters to food processing, agriculture, and assistive settings, where
produce tears, cracks, and bruises under modest compression and the damage is
irreversible~\cite{zhu2025deformable,blanco2024t,wang2025damageless}. A successful grasp must be
right in two respects at once: where and in what orientation the fingers close, and how far they
close. Closing across a thin edge, or straddling a stem, fails at a width that would have been
correct elsewhere on the same item. Both quantities depend on the item's pose, size, and shape,
and these vary continuously within a single food category and change entirely across categories,
so a policy that handles one item well need not handle its neighbour.

Existing approaches (see Section~\ref{sec:related}) address this largely through hardware and
sensing. Compliant and pneumatic end-effectors bound contact pressure
mechanically~\cite{wang2024softgripper}, and tactile or force-controlled grippers close the loop
on contact at run time, handling delicate items well~\cite{kang2026learning,vtam}. While
effective, these methods change what the robot can \emph{observe} rather than where its
\emph{supervision} comes from: the policies are learned by imitation from human demonstrations
recorded on real hardware, one item and one pose at a time. Stress-aware grasp planning targets
the damaging quantity directly~\cite{pan2020stressmin}, and finite-element simulation can
evaluate grasp outcomes on deformable objects at scale~\cite{defgraspsim}, but neither has been
used to supply the demonstrations a policy trains on. Acquiring broad coverage of the physical
variation therefore remains a human-labour problem.

That coverage is what determines generalisation. Lin \emph{et al.} show that policy performance
follows a power law in the number of environments and objects, and that demonstrations per
object saturate quickly: diversity, not count, is what buys
generalisation~\cite{lin2025datascaling}. Precision sharpens the requirement, since for
high-precision tasks the number of demonstrations needed grows super-exponentially as the
tolerance tightens, against a limit precision that itself depends on the quality of the
expert~\cite{curseofprecision}. We observe the shortfall directly on hardware: a policy trained
on 50 real demonstrations does not degrade uniformly but fails in a structured way, misgrasping
at particular workspace locations and mis-squeezing at particular object sizes
(Fig.~\ref{fig:coverage}). Data-generation systems already multiply a few hundred human
demonstrations into tens of thousands by adapting them to new object
placements~\cite{mandlekar2023mimicgen}. However, such adaptation is geometric and reuses the
grips a human chose; it does not decide afresh how firmly to close on an item it has not seen.

To address these challenges, we propose \method, a framework that generates gentle
demonstrations autonomously in simulation and trains from them a single policy spanning many
fragile food categories. Our key observation is that internal stress, the quantity that
determines whether an item is damaged and which no sensor on our robot reports, is computed
exactly inside a deformable simulator; we therefore treat gentleness as a planning objective
rather than a sensed signal. A finite-element model predicts the stress that closing to a
commanded width induces, and a scripted demonstrator grasps each item at the pose and width
that minimise it, subject to holding the item against gravity. We run this demonstrator across
\ncats{} fragile food categories while randomising object pose, size, shape, and material within
each, and collect \ndemos{} demonstrations with no human teleoperation. A single
category-conditioned policy is trained on the pooled data, co-trained with \nreal{} real
teleoperated episodes~\cite{simrealcotrain}, and deployed on a $7$-DoF arm that observes the
scene as a point cloud from one external depth camera---chosen for its narrow sim-to-real
appearance gap---with no tactile or force sensing in the loop. On hardware it grasps \ncats{}
categories at \headline{} success while handling in-category variation in pose, size, and shape.

The main contributions of this work include: 1) a stress-aware grasp planner built on a
width-controlled, distributed-pad contact model, which selects both where to grasp and how far
to close on a deformable item, 2) an autonomous demonstration generator that produces gentle,
successful grasps across randomised object pose, size, shape, and material with no human
demonstration, 3) a single category-conditioned policy that handles \ncats{} fragile food
categories from a point cloud alone, without tactile or force sensing, and 4) a real-robot
evaluation of that policy across categories and across in-category variation in pose, size, and
shape, a step towards general gentle manipulation of everyday food.
